\section{Ikhtisar Materi per Bab}

Struktur kurikulum kriptografi ini dirancang secara sistematis untuk membangun pemahaman dari konsep dasar hingga aplikasi lanjutan. Setiap bab mengintegrasikan teori matematis, implementasi praktis, dan analisis keamanan untuk memberikan pembelajaran yang komprehensif dan berorientasi aplikasi.

\begin{enumerate}
  \item \textbf{Bab II}: Pengantar kriptografi, model keamanan, ancaman, confidentiality, integrity, authenticity, kriptografi simetris vs asimetris
  \item \textbf{Bab III}: Aritmetika modular, algoritma Euclid, teori bilangan, fast exponentiation, implementasi C
  \item \textbf{Bab IV}: Cipher Caesar, substitusi monoalfabetik, Vigenère, kelemahan cipher klasik, implementasi C
  \item \textbf{Bab V}: Block cipher, DES, AES, mode operasi ECB, CBC, CFB, OFB
  \item \textbf{Bab VI}: Kriptografi asimetris, RSA, pembangkitan kunci, enkripsi/dekripsi, implementasi C
  \item \textbf{Bab VII}: Fungsi hash (SHA, MD5), integritas data, tanda tangan digital
  \item \textbf{Bab VIII}: Protokol Diffie-Hellman, TLS dasar, manajemen kunci
  \item \textbf{Bab IX}: Kriptanalisis, frequency analysis, Kasiski examination
  \item \textbf{Bab X}: Evaluasi komprehensif dan rubrik penilaian. Tinjauan pencapaian serta rekomendasi tindak lanjut.
\end{enumerate}

Perjalanan pembelajaran dimulai dengan fondasi teoretis yang kuat pada Bab II dan III, memberikan pemahaman tentang prinsip-prinsip dasar keamanan informasi dan alat matematika yang diperlukan. Bab IV dan V memperkenalkan implementasi praktis mulai dari cipher klasik hingga block cipher modern, membangun keterampilan pemrograman dan analisis secara bertahap. Pendekatan ini memastikan mahasiswa memiliki dasar yang kokoh sebelum beralih ke konsep yang lebih kompleks.

Bab VI hingga VIII membahas teknologi kriptografi modern yang secara luas digunakan dalam industri, termasuk RSA, fungsi hash, dan protokol pertukaran kunci. Materi-materi ini dirancang untuk menghubungkan konsep akademis dengan aplikasi dunia nyata yang ditemui dalam sistem keamanan digital. Bab IX menambahkan perspektif keamanan dengan mempelajari teknik kriptanalisis, memberikan pemahaman holistik tentang kekuatan dan kelemahan sistem enkripsi.

Bab penutup (X) mengintegrasikan semua pengetahuan yang telah dipelajari melalui evaluasi komprehensif dan refleksi diri. Bab ini dirancang untuk mengukur pencapaian kompetensi secara keseluruhan dan memberikan rekomendasi pengembangan lebih lanjut bagi mahasiswa. Struktur kurikulum ini memastikan lulusan memiliki pemahaman mendalam tentang kriptografi dari perspektif teoretis, praktis, dan keamanan, mempersiapkan mereka untuk karir di bidang keamanan siber.
